\documentclass[12pt,a4paper]{article}
\usepackage[brazil]{babel}
\usepackage[utf8]{inputenc}
\usepackage{amsmath, amssymb, enumitem, float, booktabs, multicol}
\usepackage{newtxtext, newtxmath}
\usepackage{graphicx}
\usepackage[hidelinks]{hyperref}
\usepackage{fancyhdr}
\usepackage{geometry, chngpage}

\usepackage{titling}

\geometry{left=3cm,right=2cm,top=3cm,bottom=2cm}

\setlength{\parskip}{0.25cm}

\pagestyle{fancy}
\fancyhf{}
\fancyhead[L]{\LaTeX: Guia do Principiante}
\fancyhead[R]{\thepage}
\fancyfoot[C]{Curso Introdutório}

\pretitle{%
	\begin{center}
		\vspace*{0.35\textheight} 
		\Huge\bfseries\color{blue}
	}
	\posttitle{\par\end{center}\vskip 1em}

\preauthor{\begin{center}\Large}
	\postauthor{\end{center}}

\predate{\begin{center}\large}
	\postdate{\end{center}}

\title{\LaTeX: Guia do Principiante}
\author{Prof. Esp. Samuel de Sousa Apolinário}
\date{Outubro de 2025}

\newcommand{\paginaautor}{
	\newpage
	\thispagestyle{empty}
	\begin{flushright}
		
		{\Large \textbf{Autor}}\\[1cm]
		\textbf{Samuel de Sousa Apolinário}\\[0.3cm]
		Licenciado em Matemática - UNEAL\\
		Especialista em Estatística Aplicada - FFOCUS\\[0.5cm]
		Contato: samukamath@outlook.com
	\end{flushright}
	\newpage
}

\begin{document}
	
	\maketitle
	\thispagestyle{empty}
	
	\newpage
	
	\newcommand{\versao}{1.0.2}

\begin{center}
	\vspace*{.3\textheight}
	
	{\Huge \color{blue}\bf \LaTeX: Guia do Iniciante}
	
	\vspace*{1cm}
	
	Versão \versao
	
	\vfill
	
	Prof. Esp. Samuel de Sousa Apolinário
	
	\vspace{15pt}
	
	Outubro de 2025
\end{center}
	\thispagestyle{empty}
	
	\newpage
	
	\paginaautor
	
	
	\tableofcontents
	\thispagestyle{empty}
	
	\newpage
	
	\section*{Prefácio}

\addcontentsline{toc}{section}{Prefácio}

O LaTeX: Guia do Iniciante nasce com o propósito de introduzir estudantes ao universo da editoração científica e técnica com o uso do LaTeX — uma poderosa ferramenta para a produção de documentos de alta qualidade tipográfica. Diferente dos editores de texto convencionais, o LaTeX prioriza a estrutura e o conteúdo, permitindo que o autor concentre seus esforços na clareza das ideias, enquanto o sistema se encarrega da formatação e da estética do texto.

Esta apostila foi elaborada com uma linguagem acessível, voltada especialmente para quem está dando os primeiros passos. O conteúdo é apresentado de forma gradual. Cada seção busca aliar teoria e prática, incentivando o leitor a experimentar os exemplos propostos e a adaptar os códigos para suas próprias necessidades.

Mais do que um manual técnico, este guia pretende ser um companheiro de aprendizado, despertando o interesse pela organização e pela beleza que o LaTeX proporciona na escrita científica. Ao longo das páginas, o leitor perceberá que dominar o LaTeX não é apenas aprender comandos, mas compreender uma nova forma de pensar a produção de textos — com precisão, elegância e profissionalismo.

Espero que este material contribua para tornar sua jornada no LaTeX mais leve, instigante e produtiva.
	\thispagestyle{empty}
	
	\newpage
	
	% =============================
	\section{Qual a finalidade do LaTeX?}
	O LaTeX é uma ferramenta para formatação e elaboração de documentos voltada à criação de textos científicos, técnicos e acadêmicos com excelente qualidade tipográfica. É amplamente empregado em campos como Matemática, Engenharia, Física, Computação e demais áreas das Ciências.
	
	
	\textbf{Características principais:}
	\begin{itemize}
		\item \textbf{Qualidade tipográfica excepcional}
		
		O LaTeX produz textos com uma diagramação impecável, especialmente para fórmulas matemáticas, tabelas e citações. O resultado final parece de editora profissional;
		
		\item \textbf{Separação entre conteúdo e formatação}
		
		Você escreve o conteúdo (texto, equações, figuras) e o LaTeX cuida automaticamente da aparência. Isso evita perda de tempo com ajustes manuais de layout;
		
		\item \textbf{Excelente suporte a matemática}
		
		Nenhum editor de texto comum chega perto do poder do LaTeX para expressões matemáticas — ele foi literalmente criado com isso em mente.
		
		\item \textbf{Ideal para textos científicos e acadêmicos}
		
		Ele organiza referências, citações, índices, sumários e bibliografias automaticamente.
		Pacotes como \textbf{biblatex} ou \textbf{abntex2} (no Brasil) facilitam a vida acadêmica.
		
		\item Base em código
		
		Tudo é feito por comandos de texto, o que dá controle total sobre o documento e facilita versionamento com Git (ótimo pra quem trabalha em equipe).
		
		\item \textbf{Altamente personalizável}
		
		Existem milhares de pacotes (extensões) pra quase tudo: gráficos, tabelas, diagramas, apresentações (como o Beamer) e até currículos.
		
		\item \textbf{Compatibilidade e portabilidade}
		
		Os arquivos .tex são leves e funcionam em qualquer sistema operacional. O resultado final é geralmente um PDF padronizado.
		
		\item \textbf{Estabilidade e precisão}
		
		Mesmo documentos enormes (como teses ou livros) mantêm consistência no estilo e numeração automática impecável — sem os “bugs” comuns de editores WYSIWYG.
	
	\end{itemize}
	
	% =============================
	\section{Instalação do LaTeX}
	
	\subsection{Opção 1: Usar o Overleaf (online)}
	A forma mais simples para começar é usando o site:
	\url{https://www.overleaf.com}
	
	Não precisa instalar nada. Basta criar uma conta e começar um novo projeto.
	
	\subsection{Opção 2: Instalação local}
	Instale um compilador LaTeX e um editor:
	
	\begin{itemize}
		\item Windows: instalar o \textbf{MiKTeX} + \textbf{TeXstudio};
		\item Linux: instalar o pacote \textbf{texlive-full} + \textbf{TeXstudio};
		\item MacOS: instalar o \textbf{MacTeX}.
	\end{itemize}
	
	% =============================
	\section{Primeiro Documento em LaTeX}
	
	Vamos escrever nosso primeiro código:
	
	\begin{verbatim}
		\documentclass{article}
		\begin{document}
			Olá, mundo!
		\end{document}
	\end{verbatim}
	
	\textbf{Explicação:}
	
	\begin{itemize}
		\item \verb|\documentclass{article}|: esta linha especifica o tipo de documento que estamos produzindo. No caso, escolhemos “article”, que é indicado para textos curtos, relatórios ou artigos acadêmicos.
		
		\item \verb|\begin{document}| e \verb|\end{document}|: tudo o que estiver entre essas instruções será considerado conteúdo do documento, ou seja, o que aparecerá no PDF final.
	\end{itemize}
	
	\textbf{Como funciona}:
	
	Quando você compilar este código, o \LaTeX\, processará as instruções e gerará um arquivo PDF contendo apenas a frase “Olá, mundo!”. Essa é a forma mais básica de testar se o ambiente está funcionando corretamente e de entender a estrutura mínima de um documento \LaTeX.
	
	Observações importantes:
	
	\begin{itemize}
		
		\item Todo documento em \LaTeX precisa ter uma classe definida (\verb|\documentclass|), que indica o formato e estilo geral do arquivo.
		
		\item O conteúdo visível sempre deve estar entre \verb|\begin{document}| e \verb|\end{document}|. Qualquer texto colocado fora dessas linhas não aparecerá no documento final.
		
		\item Esse exemplo é chamado de documento mínimo ou exemplo inicial, útil para compreender a organização básica de um arquivo \LaTeX antes de adicionar elementos mais complexos, como títulos, listas ou imagens.
		
	\end{itemize}
	
	Dica: experimente alterar o texto “Olá, mundo!” para ver como o \LaTeX reproduz o conteúdo que você escreve, mantendo a formatação automática do documento.
	% =============================
	\section{Formatando Textos}
	
	O \LaTeX{} oferece diversos recursos para formatar textos, permitindo destacar informações importantes e organizar melhor o conteúdo do documento. A seguir, veremos alguns dos comandos mais usados na formatação básica:
	
	\begin{itemize}
		\item \textbf{Negrito:} utilizado para dar ênfase a palavras ou expressões importantes. O comando é feito com \verb|\textbf{texto}| e o resultado é \textbf{texto em negrito}.
		\item \textit{Itálico:} serve para destacar termos estrangeiros, títulos de obras ou dar ênfase leve. Use \verb|\textit{texto}| e o resultado é \textit{texto em itálico}.
		\item \underline{Sublinhado:} pode ser aplicado com \verb|\underline{texto}|, gerando um \underline{texto sublinhado}. No entanto, é recomendado usar com moderação, já que o sublinhado é menos comum em textos acadêmicos.
	\end{itemize}
	
	Além dessas opções, o \LaTeX{} permite controlar a estrutura dos parágrafos e a quebra de linhas de forma simples:
	
	\begin{itemize}
		\item Para pular linhas dentro de um mesmo parágrafo, utilize o comando \verb|\\|.  
		\textbf{Exemplo:}  
		\verb|Primeira linha.\\Segunda linha.| produz:  
		Primeira linha.\\Segunda linha.
		\item Para iniciar um novo parágrafo, basta deixar uma linha em branco no código. O \LaTeX{} automaticamente adiciona o espaçamento adequado entre parágrafos.
	\end{itemize}
	
	
	
	É importante lembrar que o \LaTeX{} foi projetado para cuidar automaticamente da formatação do texto, garantindo uma apresentação limpa e uniforme. Por isso, é sempre bom evitar o uso excessivo de comandos de estilo e deixar que o sistema organize o layout do documento.
	\\ 
	
	Também é possível combinar estilos, como \verb|\textbf{\textit{texto}}|, para criar um \textbf{\textit{texto em negrito e itálico}}. Essa combinação é útil quando se deseja dar destaque ainda maior a um termo específico.
	
	\subsection{Tamanho da fonte}
	
	O \LaTeX{} permite ajustar o tamanho das letras de forma rápida e padronizada, utilizando comandos específicos que alteram a formatação apenas do trecho desejado. Essa flexibilidade é útil para destacar títulos, legendas, citações, notas de rodapé e outros elementos do texto.
	
	\begin{itemize}
		\item \verb|\tiny| — 5pt
		\item \verb|\scriptsize| — 7pt
		\item \verb|\footnotesize| — 8pt
		\item \verb|\small| — 9pt
		\item \verb|\normalsize| — 10pt (tamanho padrão)
		\item \verb|\large| — 12pt
		\item \verb|\Large| — 14.4pt
		\item \verb|\LARGE| — 17.28pt
		\item \verb|\huge| — 20.74pt
		\item \verb|\Huge| — 24.88pt
	\end{itemize}
	
	Esses comandos podem ser aplicados a qualquer parte do texto. Por exemplo:
	
	\begin{verbatim}
		{\large Este texto está um pouco maior.}
		
		{\Huge Este texto está bem grande!}
	\end{verbatim}
	
	O resultado será:
	
	{\large Este texto está um pouco maior.}
	
	{\Huge Este texto está bem grande!}
	
	\bigskip
	Note que as chaves \verb|{ }| delimitam o trecho do texto que será afetado. Assim, o tamanho da fonte volta ao normal após o fechamento das chaves.  
	Se você quiser alterar o tamanho de uma seção inteira, pode colocar o comando antes do texto, sem chaves, por exemplo:
	
	\begin{verbatim}
		\huge
		Este parágrafo inteiro ficará em tamanho grande até
		que outro comando de tamanho seja usado.
	\end{verbatim}
	
	\bigskip
	\textbf{Dica:} é uma boa prática usar variações de tamanho com moderação — apenas para destacar partes importantes, como títulos, seções ou observações especiais. O excesso pode deixar o texto visualmente cansativo.
	
	\bigskip
	Também é possível definir o tamanho base do documento ao criar o arquivo \LaTeX{}, adicionando um parâmetro à classe do documento. Por exemplo:
	
	\begin{verbatim}
		\documentclass[12pt]{article}
	\end{verbatim}
	
	Nesse caso, o tamanho padrão (\verb|\normalsize|) será 12 pontos, e os demais tamanhos serão ajustados proporcionalmente.  
	Esse recurso é bastante útil para atender a exigências de formatação de universidades, revistas científicas ou editais específicos.
	
	\bigskip
	\textbf{Resumo:}  
	Os tamanhos de fonte no \LaTeX{} seguem uma hierarquia lógica — do menor para o maior — e garantem uma aparência harmônica em todo o documento, mantendo o padrão tipográfico profissional característico do sistema.
	
	% =============================
	\section{Listas}
	
	O \LaTeX{} permite criar diferentes tipos de listas para organizar informações de forma clara e estruturada. As listas são recursos muito úteis em textos acadêmicos, relatórios e materiais didáticos, pois ajudam na leitura e na hierarquização de ideias.
	
	\subsection{Lista não numerada}
	
	As listas não numeradas são usadas quando a ordem dos itens não é relevante. São criadas com o ambiente \verb|itemize|. Veja o exemplo:
	
	\begin{figure}[H]
		\begin{minipage}{.45\linewidth}
			\begin{verbatim}
				\begin{itemize}
					\item Primeiro item
					\item Segundo item
					\item Terceiro item
				\end{itemize}
			\end{verbatim}
		\end{minipage}
		\begin{minipage}{.45\linewidth}
			O resultado será exibido assim:
			
			\begin{itemize}
				\item Primeiro item
				\item Segundo item
				\item Terceiro item
			\end{itemize}
			
		\end{minipage}
	\end{figure}
	
	Por padrão, o \LaTeX{} utiliza marcadores simples (como bolinhas) antes de cada item. É possível criar listas aninhadas, ou seja, uma lista dentro de outra, para mostrar subitens relacionados. Exemplo:
	
	\begin{figure}[H]
		\begin{minipage}{.45\linewidth}
			
			\begin{verbatim}
				\begin{itemize}
					\item Frutas
					\begin{itemize}
						\item Maçã
						\item Banana
						\item Laranja
					\end{itemize}
					\item Verduras
					\item Legumes
				\end{itemize}
			\end{verbatim}
		\end{minipage}
		\begin{minipage}{.45\linewidth}
			Que será exibido assim:
			
			\begin{itemize}
				\item Frutas
				\begin{itemize}
					\item Maçã
					\item Banana
					\item Laranja
				\end{itemize}
				\item Verduras
				\item Legumes
			\end{itemize}
		\end{minipage}
	\end{figure}
	
	Esse tipo de estrutura facilita a organização de informações em tópicos e sub-tópicos, tornando o texto mais visual e compreensível.
	
	Também é possível personalizar os símbolos das listas com o pacote \verb|enumitem|, por exemplo:
	\begin{verbatim}
		\usepackage{enumitem}
		
		\begin{itemize}[label=--]
			\item Item com traço
			\item Outro item
		\end{itemize}
	\end{verbatim}
	
	Isso substitui as bolinhas por traços, o que pode deixar o visual mais limpo em apresentações ou relatórios.
	
	\bigskip
	Outros tipos de listas incluem:
	\begin{itemize}
		\item \textbf{Listas numeradas} – criadas com o ambiente \verb|enumerate|, usadas quando a sequência dos itens importa;
		\item \textbf{Listas de descrição} – feitas com o ambiente \verb|description|, ideais para definir termos ou conceitos.
	\end{itemize}
	
	\subsection{Lista numerada}
	
	As listas numeradas são ideais quando a ordem dos itens tem importância, como em etapas de um procedimento, instruções ou classificações. No \LaTeX{}, elas são criadas com o ambiente \verb|enumerate|. Veja o exemplo básico:
	
	\begin{figure}[H]
		\begin{minipage}{.45\linewidth}
			\begin{verbatim}
				\begin{enumerate}
					\item Primeiro
					\item Segundo
					\item Terceiro
				\end{enumerate}
			\end{verbatim}
		\end{minipage}
		\begin{minipage}{.45\linewidth}
			O resultado será exibido assim:
			
			\begin{enumerate}
				\item Primeiro
				\item Segundo
				\item Terceiro
			\end{enumerate}
		\end{minipage}
	\end{figure}
	
	Cada item é numerado automaticamente pelo \LaTeX{}, o que significa que, mesmo que você adicione ou remova itens, a numeração será ajustada de forma automática — uma das grandes vantagens em relação à digitação manual.
	
	Também é possível criar listas numeradas dentro de outras listas (listas aninhadas), permitindo organizar informações hierarquicamente. Veja o exemplo:
	
	\begin{figure}[H]
		\begin{minipage}{.45\linewidth}
			\begin{verbatim}
				\begin{enumerate}
					\item Etapa 1
					\begin{enumerate}
						\item Subetapa 1.1
						\item Subetapa 1.2
					\end{enumerate}
					\item Etapa 2
				\end{enumerate}
			\end{verbatim}
		\end{minipage}
		\begin{minipage}{.45\linewidth}
			E o resultado será:
			
			\begin{enumerate}
				\item Etapa 1
				\begin{enumerate}
					\item Subetapa 1.1
					\item Subetapa 1.2
				\end{enumerate}
				\item Etapa 2
			\end{enumerate}
		\end{minipage}
	\end{figure}
	
	Por padrão, o \LaTeX{} usa números arábicos (1, 2, 3, ...), mas é possível personalizar o formato da numeração com o pacote \verb|enumitem|. Por exemplo:
	
	\begin{verbatim}
		\usepackage{enumitem}
		
		\begin{enumerate}[label=\alph*)]
			\item Primeira opção
			\item Segunda opção
		\end{enumerate}
	\end{verbatim}
	
	Esse comando faz com que os itens sejam listados com letras minúsculas seguidas de parêntese, resultando em:
	\begin{enumerate}[label=\alph*)]
		\item Primeira opção
		\item Segunda opção
	\end{enumerate}
	
	Outros formatos possíveis incluem:
	\begin{itemize}
		\item \verb|label=\roman*)| → numeração com algarismos romanos minúsculos (i, ii, iii, ...);
		\item \verb|label=\Roman*)| → algarismos romanos maiúsculos (I, II, III, ...);
		\item \verb|label=\Alph*)| → letras maiúsculas (A, B, C, ...).
	\end{itemize}
	
	Essas personalizações permitem ajustar o estilo das listas conforme o tipo de documento, mantendo a formatação elegante e consistente.
	
	% =============================
	\section{Inserindo Títulos e Sumário}
	
	O \LaTeX{} facilita muito a criação automática de títulos, autores, datas e até sumários, sem a necessidade de formatação manual. Esses recursos são especialmente úteis em trabalhos acadêmicos, relatórios e livros, pois mantêm a padronização e permitem atualizações automáticas sempre que o conteúdo é alterado.
	
	\begin{verbatim}
		\title{Meu Documento}
		\author{Nome do Autor}
		\date{Data}
		\maketitle
		\tableofcontents
		\section{Introdução}
	\end{verbatim}
	
	Esses comandos geram automaticamente uma \textbf{capa simples} e um \textbf{sumário}. A seguir, entenda o papel de cada um:
	
	\begin{itemize}
		\item \verb|\title{...}| define o título principal do documento. O texto colocado entre as chaves será exibido em destaque no topo da capa.
		\item \verb|\author{...}| define o(s) nome(s) do(s) autor(es) do trabalho. É possível incluir mais de um autor separando-os com barras verticais (\verb|\and|).
		\item \verb|\date{...}| insere a data. Você pode digitar uma data específica, como \verb|\date{Outubro de 2025}|, ou deixar o campo vazio \verb|\date{}| para omitir a data. Se desejar inserir a data atual automaticamente, use \verb|\date{\today}|.
		\item \verb|\maketitle| gera a página de título com as informações fornecidas acima.
		\item \verb|\tableofcontents| cria automaticamente o \textbf{sumário}, com base nos títulos e subtítulos definidos com \verb|\section|, \verb|\subsection|, etc.
	\end{itemize}
	
	Por exemplo, o código:
	
	\begin{verbatim}
		\documentclass{article}
		\title{Relatório de Experimentos}
		\author{Maria Silva e João Pereira}
		\date{\today}
		
		\begin{document}
			\maketitle
			\tableofcontents
			
			\section{Introdução}
			\section{Metodologia}
			\section{Resultados}
			\section{Conclusão}
		\end{document}
	\end{verbatim}
	
	gera um documento com capa, data atual automática e um sumário estruturado que se atualiza conforme você adiciona ou remove seções.
	
	\bigskip
	\textbf{Dica:} se o sumário não estiver aparecendo atualizado, compile o documento duas vezes. O \LaTeX{} precisa de uma segunda execução para gerar corretamente o índice, já que ele é montado a partir das informações coletadas durante a primeira compilação.
	
	\bigskip
	Você também pode personalizar o título e o sumário com pacotes adicionais, como:
	\begin{itemize}
		\item \verb|titlesec| — permite mudar o estilo e o espaçamento dos títulos;
		\item \verb|tocloft| — possibilita modificar a aparência do sumário (margens, fonte, espaçamento, etc.).
	\end{itemize}

	
	% =============================
	\section{Equações Matemáticas}
	
	O \LaTeX{} é amplamente reconhecido por sua capacidade de escrever expressões matemáticas com alta qualidade tipográfica. Ele utiliza um modo especial de escrita chamado \textbf{modo matemático}, que permite representar símbolos, frações, expoentes, raízes, somatórios, integrais e muito mais.
	
	Existem basicamente duas formas principais de inserir equações: o \textbf{modo inline} e o \textbf{modo centralizado}. Além disso, é possível criar sistemas de equações e expressões alinhadas.
	
	\subsection{Modo inline}
	
	O modo \textit{inline} é usado quando a equação aparece no meio de um parágrafo, acompanhando o texto normalmente.  
	Por exemplo:
	
	\begin{verbatim}
		A equação é $E = mc^2$.
	\end{verbatim}
	
	O resultado será:  
	A equação é $E = mc^2$.
	
	Note que a equação é inserida entre os símbolos de cifrão (\verb|$...$|). Essa é a forma mais simples e comum de escrever expressões curtas, como $a^2 + b^2 = c^2$ ou $\pi \approx 3.14$.
	
	\bigskip
	\textbf{Dica:} use o modo inline apenas para fórmulas pequenas. Fórmulas longas ou que precisem de destaque devem ser centralizadas.
	
	\subsection{Modo centralizado}
	
	Quando a expressão matemática é o foco da explicação, o ideal é centralizá-la no documento. Para isso, usamos o ambiente delimitado por \verb|\[| e \verb|\]|:
	
	\begin{verbatim}
		\[
		x = \frac{-b \pm \sqrt{b^2 - 4ac}}{2a}
		\]
	\end{verbatim}
	
	O resultado será:
	
	\[
	x = \frac{-b \pm \sqrt{b^2 - 4ac}}{2a}
	\]
	
	Esse formato é ideal para equações destacadas, demonstrações ou passos importantes de um cálculo.
	
	\bigskip
	Também é possível numerar as equações automaticamente utilizando o ambiente \verb|equation|:
	
	\begin{verbatim}
		\begin{equation}
			E = mc^2
		\end{equation}
	\end{verbatim}
	
	O resultado será:
	
	\begin{equation}
		E = mc^2
	\end{equation}
	
	O número gerado pode ser usado para referência cruzada no texto com o comando \verb|\label| e \verb|\ref|, por exemplo:
	\begin{verbatim}
		\begin{equation}\label{eq:energia}
			E = mc^2
		\end{equation}
		A equação \ref{eq:energia} mostra a equivalência entre massa e energia.
	\end{verbatim}
	
	\subsection{Sistema de equações}
	
	Para representar um conjunto de equações relacionadas, podemos usar o ambiente \verb|cases|, dentro do modo matemático centralizado.  
	Exemplo:
	
	\begin{verbatim}
		\[
		\begin{cases}
			x + y = 10 \\
			2x - y = 3
		\end{cases}
		\]
	\end{verbatim}
	
	O resultado será:
	
	\[
	\begin{cases}
		x + y = 10 \\
		2x - y = 3
	\end{cases}
	\]
	
	Esse tipo de estrutura é especialmente útil para representar sistemas lineares, condições por partes e definições de funções.
	
	\bigskip
	\textbf{Dica extra:} para equações mais elaboradas, com alinhamentos e múltiplas linhas, o pacote \verb|amsmath| oferece ambientes como:
	\begin{itemize}
		\item \verb|align| — alinha várias equações pelo sinal de igualdade;
		\item \verb|gather| — centraliza várias equações sem alinhamento específico;
		\item \verb|split| — divide uma equação longa em várias linhas.
	\end{itemize}
	
	Esses ambientes deixam o texto matemático mais organizado e legível, além de garantirem formatação profissional.
	
	
	% =============================
	\section{Inserindo Imagens}
	
	O \LaTeX{} permite inserir imagens de maneira simples e profissional, garantindo que elas fiquem bem posicionadas e integradas ao texto. Para isso, utilizamos o pacote \verb|graphicx|, que oferece o comando \verb|\includegraphics| e diversas opções de personalização.
	
	Antes de tudo, é necessário carregar o pacote no início do documento, dentro do preâmbulo:
	
	\begin{verbatim}
		\usepackage{graphicx}
	\end{verbatim}
	
	Depois disso, podemos inserir uma imagem utilizando o ambiente \verb|figure|, como no exemplo abaixo:
	
	\begin{verbatim}
		\begin{figure}[h]
			\centering
			\includegraphics[width=0.5\textwidth]{exemplo.jpg}
			\caption{Minha imagem}
		\end{figure}
	\end{verbatim}
	
	Esse código gera uma figura centralizada, com legenda e tamanho ajustado proporcionalmente à largura do texto.
	
	\bigskip
	\textbf{Explicação dos comandos:}
	
	\begin{itemize}
		\item \verb|\begin{figure}[h]| — inicia o ambiente de figura.  
			O argumento opcional \verb|[h]| sugere ao \LaTeX{} que a imagem seja posicionada “aqui” (*here*), próximo ao local onde foi inserida no código.  
			Outras opções possíveis são:
			\begin{itemize}
				\item \verb|t| – posiciona a imagem no topo da página (*top*);
				\item \verb|b| – posiciona no rodapé da página (*bottom*);
				\item \verb|p| – coloca a figura em uma página exclusiva para imagens (*page*);
				\item \verb|!| – força o \LaTeX{} a tentar obedecer o local indicado.
			\end{itemize}
			
			\item \verb|\centering| — centraliza a imagem dentro do ambiente.
			
			\item \verb|\includegraphics[width=0.5\textwidth]{exemplo.jpg}| — insere a imagem “exemplo.jpg” ajustada a 50\% da largura do texto.  
			Você também pode usar:
			\begin{itemize}
				\item \verb|width| — define a largura (ex.: \verb|width=8cm|);
				\item \verb|height| — define a altura (ex.: \verb|height=4cm|);
				\item \verb|scale| — redimensiona proporcionalmente (ex.: \verb|scale=0.8|).
			\end{itemize}
			
			\item \verb|\caption{Minha imagem}| — adiciona uma legenda abaixo da figura. O \LaTeX{} numera automaticamente as figuras conforme a ordem em que aparecem.
			
			\item \verb|\end{figure}| — finaliza o ambiente da figura.
	\end{itemize}
	
	\bigskip
	\textbf{Exemplo completo:}
	
	\begin{verbatim}
		\documentclass{article}
		\usepackage{graphicx}
		\begin{document}
			
			\begin{figure}[h]
				\centering
				\includegraphics[width=0.6\textwidth]{paisagem.png}
				\caption{Vista panorâmica da montanha}
				\label{fig:montanha}
			\end{figure}
			
			A Figura~\ref{fig:montanha} mostra uma bela vista das montanhas 
			durante o amanhecer.
		\end{document}
	\end{verbatim}
	
	\bigskip
	\textbf{Dicas importantes:}
	\begin{itemize}
		\item As imagens devem estar na mesma pasta do arquivo \verb|.tex| ou em um diretório especificado (ex.: \verb|\includegraphics{imagens/exemplo.jpg}|).
		\item É recomendável usar formatos como \texttt{.png}, \texttt{.jpg} ou \texttt{.pdf} (este último é ideal para gráficos vetoriais).
		\item Você pode combinar legendas e rótulos para fazer referência cruzada a figuras no texto.
		\item Para colocar duas imagens lado a lado, utilize o pacote \verb|subfigure| ou \verb|subcaption|.
	\end{itemize}
	
	Essas opções tornam a inserção de imagens no \LaTeX{} muito flexível, mantendo sempre o padrão de qualidade tipográfica que é a marca registrada do sistema.

	% =============================
	\section{Tabelas Simples}
	
	O \LaTeX{} também permite criar tabelas de maneira estruturada e precisa, garantindo alinhamento perfeito e aparência profissional.  
	O ambiente básico usado para isso é o \verb|tabular|, que define as colunas e o conteúdo de cada célula da tabela.
	
	\begin{verbatim}
		\begin{tabular}{|c|c|}
			\hline
			Nome & Nota \\
			\hline
			João & 8,5 \\
			Ana & 9,0 \\
			\hline
		\end{tabular}
	\end{verbatim}
	
	Esse código gera uma tabela simples, com bordas verticais e horizontais, da seguinte forma:
	
	\begin{center}
		\begin{tabular}{|c|c|}
			\hline
			Nome & Nota \\
			\hline
			João & 8,5 \\
			Ana & 9,0 \\
			\hline
		\end{tabular}
	\end{center}
	
	\bigskip
	\textbf{Explicando cada parte:}
	
	\begin{itemize}
		\item \verb|\begin{tabular}{|\textbar \verb|c| \textbar \verb|c| \textbar \verb|}| — inicia o ambiente da tabela e define o formato das colunas.  
			O argumento entre chaves indica:
			\begin{itemize}
				\item \verb|c| — centraliza o conteúdo da coluna;  
				\item \verb|l| — alinha o conteúdo à esquerda (*left*);  
				\item \verb|r| — alinha o conteúdo à direita (*right*);  
				\item \verb"| |" — adiciona linhas verticais entre as colunas.
			\end{itemize}
			
			\item \verb|\hline| — insere uma linha horizontal acima ou abaixo das células.
			
			\item O símbolo \verb|&| separa as colunas dentro de uma linha.
			
			\item O comando \verb|\\| indica o fim de uma linha e o início da próxima.
		\end{itemize}
		
		\bigskip
		\textbf{Exemplo com três colunas:}
		
		\begin{verbatim}
			\begin{tabular}{|l|c|r|}
				\hline
				Aluno & Disciplina & Nota \\
				\hline
				João & Matemática & 8,5 \\
				Ana & Português & 9,0 \\
				Marcos & História & 7,8 \\
				\hline
			\end{tabular}
		\end{verbatim}
		
		Resultado:
		
		\begin{center}
			\begin{tabular}{|l|c|r|}
				\hline
				Aluno & Disciplina & Nota \\
				\hline
				João & Matemática & 8,5 \\
				Ana & Português & 9,0 \\
				Marcos & História & 7,8 \\
				\hline
			\end{tabular}
		\end{center}
		
		\bigskip
		\textbf{Dicas úteis:}
		\begin{itemize}
			\item Evite excesso de linhas verticais e horizontais — tabelas muito marcadas visualmente podem ficar poluídas.
			\item Use o pacote \verb|booktabs| para criar tabelas com estilo mais limpo e elegante (usando comandos como \verb|\toprule|, \verb|\midrule| e \verb|\bottomrule|).
			\item Para adicionar título e numeração automática à tabela, utilize o ambiente \verb|table| com o comando \verb|\caption|.
		\end{itemize}
		
		\bigskip
		\textbf{Exemplo com legenda:}
		
		\begin{verbatim}
			\begin{table}[h]
				\centering
				\begin{tabular}{|c|c|}
					\hline
					Produto & Preço (R\$) \\
					\hline
					Caderno & 15,00 \\
					Caneta & 3,50 \\
					\hline
				\end{tabular}
				\caption{Tabela de preços de materiais escolares}
			\end{table}
		\end{verbatim}
		
		Esse código cria uma tabela com legenda e numeração automática, permitindo fazer referência cruzada dentro do texto.
		
	
	% =============================
	\section{Criando um Documento de Prova}
	
	O \LaTeX{} é uma excelente ferramenta para criar provas e atividades avaliativas com qualidade visual e organização. Ele permite combinar texto, expressões matemáticas e formatação consistente, o que resulta em documentos elegantes e fáceis de ler — tanto para o professor quanto para o aluno.
	
	\bigskip
	\textbf{Passos iniciais:}
	
	\begin{itemize}
		\item Use a classe de documento padrão para artigos:
		\begin{verbatim}
			\documentclass[12pt]{article}
		\end{verbatim}
		A opção \verb|12pt| define o tamanho da fonte, deixando o texto mais legível para impressões.  
		Outros tamanhos possíveis incluem \verb|10pt| e \verb|11pt|.
		
		\item Acrescente título, data e um espaço reservado para o nome do aluno:
		\begin{verbatim}
			\title{Prova de Matemática}
			\author{Professor João Silva}
			\date{Data: \underline{\hspace{3cm}} \\
				Nome: \underline{\hspace{8cm}}}
		\end{verbatim}
		Assim, o título e o campo para o nome do aluno aparecem automaticamente no cabeçalho quando usamos \verb|\maketitle|.
		
		\item Utilize comandos matemáticos para escrever as questões com clareza e precisão, aproveitando o modo matemático do \LaTeX{}:
		\begin{verbatim}
			1) Calcule o valor de $x$ na equação $2x + 5 = 11$.
			
			2) Resolva o sistema:
			\[
			\begin{cases}
				3x + y = 7 \\
				x - y = 1
			\end{cases}
			\]
		\end{verbatim}
	\end{itemize}
	
	\bigskip
	\textbf{Exemplo completo de documento:}
	
	\begin{verbatim}
		\documentclass[12pt]{article}
		\usepackage{amsmath}
		\usepackage{geometry}
		\geometry{a4paper, margin=2.5cm}
		
		\title{Prova de Matemática - 7º Ano}
		\author{Professor João Silva}
		\date{Data: \underline{\hspace{3cm}} \\
			Nome: \underline{\hspace{8cm}}}
		
		\begin{document}
			\maketitle
			
			\section*{Questões}
			
			\begin{enumerate}
				\item Resolva a equação $3x + 7 = 16$.
				
				\item Calcule a área de um triângulo de base $b = 8\,cm$ e altura $h = 5\,cm$.
				
				\item Resolva o sistema:
				\[
				\begin{cases}
					2x + y = 10 \\
					x - y = 4
				\end{cases}
				\]
				
				\item Simplifique a fração $\dfrac{18}{24}$.
			\end{enumerate}
			
		\end{document}
	\end{verbatim}
	
	\bigskip
	\textbf{Dicas úteis:}
	\begin{itemize}
		\item Use o pacote \verb|geometry| para ajustar margens conforme o padrão da escola.
		\item Adicione \verb|\newpage| para separar seções ou páginas quando a prova for longa.
		\item O pacote \verb|multicol| permite dividir as questões em colunas, útil para economizar espaço.
		\item Você pode inserir imagens, gráficos e figuras com \verb|graphicx|, se a prova incluir interpretação visual.
		\item Para criar linhas de resposta, use comandos como \verb|\underline{\hspace{5cm}}| ou o ambiente \verb|tabular|.
	\end{itemize}
	
	Com essas configurações, é possível gerar avaliações limpas, bem organizadas e de aparência profissional, economizando tempo e mantendo consistência nas futuras provas.
	
	
	% =============================
	\section{Comandos Úteis}
	
	O \LaTeX{} oferece diversos comandos práticos para controlar espaçamentos, quebras de página e outros ajustes finos no layout do documento. Esses comandos são especialmente úteis quando você quer refinar a aparência do texto sem alterar a estrutura geral do conteúdo.
	
	\begin{itemize}
		\item \verb|\newpage| \texttwelveudash\, força uma quebra de página, fazendo com que o conteúdo seguinte comece na próxima página.  
		Exemplo:
		\begin{verbatim}
			Primeira parte do texto.
			\newpage
			Segunda parte começa em uma nova página.
		\end{verbatim}
		
		\bigskip
		
		\item \verb|\vspace{1cm}| \texttwelveudash\, insere um espaço vertical de 1 cm entre linhas ou parágrafos. Pode ser usado para dar “respiro” em partes específicas do texto.  
		Exemplo:
		\begin{verbatim}
			Texto acima.
			
			\vspace{1cm}
			
			Texto abaixo.
		\end{verbatim}
		
		Resultado: um espaço vertical de 1 cm entre os dois trechos de texto.
		
		\bigskip
		
		\item \verb|\hspace{2cm}| \texttwelveudash\, adiciona um espaço horizontal de 2 cm na mesma linha.  
		Esse comando é útil, por exemplo, para criar deslocamentos, alinhar trechos manualmente ou separar elementos no mesmo parágrafo.  
		Exemplo:
		\begin{verbatim}
			Nome:\hspace{2cm}Assinatura:
		\end{verbatim}
		
		Resultado: entre \textquotedblleft Nome:\textquotedblright\, e \textquotedblleft Assinatura:\textquotedblright\, haverá um espaço de 2 cm.
		
		\bigskip
		
		\item \verb|~| \texttwelveudash\, cria um espaço \textbf{rígido} entre palavras, impedindo que o \LaTeX{} quebre a linha nesse ponto.  
		É muito usado em situações como abreviações seguidas de nomes, números e unidades de medida.  
		Exemplo:
		\begin{verbatim}
			Prof.~Silva publicou o artigo em 2025.
		\end{verbatim}
		
		Resultado: o \textquotedblleft Prof.\textquotedblright \, e o \textquotedblleft Silva\textquotedblright\, permanecem juntos, sem quebra de linha entre eles.
	\end{itemize}
	
	\bigskip
	\textbf{Dicas práticas:}
	\begin{itemize}
		\item Para inserir um espaço vertical que também seja válido no início de uma nova página, use \verb|\vspace*{...}| (note o asterisco).
		\item Para múltiplas quebras de página seguidas, prefira \verb|\clearpage| — ela garante que todas as figuras e tabelas pendentes sejam impressas antes da nova página.
		\item Combine \verb|\vspace| e \verb|\hspace| para criar áreas em branco personalizadas em certificados, formulários ou modelos de avaliação.
	\end{itemize}
	
	Esses comandos, apesar de simples, ajudam a refinar a apresentação e tornam o controle visual do documento muito mais preciso.
	
	
	% =============================
	
	\section{Caracteres especiais}
	
	Alguns símbolos e caracteres têm funções específicas dentro do \LaTeX{}, e por isso não podem ser usados diretamente no texto sem causar erros de compilação. Eles são chamados de \textbf{caracteres especiais}, pois fazem parte da sintaxe do sistema.
	
	A tabela abaixo mostra os principais deles:
	
	\begin{multicols}{5}
		\begin{itemize}
			\item \verb|$|
			\item \verb|#|
			\item \verb|%|
			\item \verb|^|
			\item \verb|_|
			\item \verb|{ }|
			\item \verb|~|
			\item \verb|\|
			\item \verb|&|
		\end{itemize}
	\end{multicols}
	
	\bigskip
	Esses caracteres são interpretados pelo compilador como comandos ou delimitadores, e não como texto comum.  
	Por exemplo, o símbolo \verb|$| é usado para abrir e fechar o modo matemático, e o \verb|%| serve para inserir comentários no código.
	
	\bigskip
	\textbf{Como exibir caracteres especiais no texto}
	
	Se quiser mostrar esses símbolos no corpo do texto, é necessário usar uma forma de “escape”, ou seja, precedê-los com uma barra invertida (\verb|\|) ou comandos específicos. Veja alguns exemplos:
	

		\begin{itemize}
			\item \verb|\$| → exibe \$
			\item \verb|\#| → exibe \#
			\item \verb|\%| → exibe \%
			\item \verb|\^{}| → exibe \^{} (acento circunflexo isolado)
			\item \verb|\_| → exibe \_
			\item \verb|\{| e \verb|\}| → exibem \{ e \}
			\item \verb|\~{}| → exibe \~{} (til isolado)
			\item \verb|\textbackslash| → exibe uma barra invertida: \textbackslash
			\item \verb|\&| → exibe \&
		\end{itemize}
	
	\bigskip
	\textbf{Função de cada um no \LaTeX{}:}
	
	\begin{itemize}
		\item \textbf{\$} — abre e fecha o modo matemático.  
		Exemplo: \verb|$a^2 + b^2 = c^2$| → $a^2 + b^2 = c^2$.
		
		\item \textbf{\#} — usado para indicar parâmetros dentro de comandos e macros personalizados.
		
		\item \textbf{\%} — inicia um comentário. Tudo o que vem após o símbolo, até o fim da linha, é ignorado pelo compilador.
		
		\item \textbf{\^} e \textbf{\_} — utilizados em expressões matemáticas para sobrescritos (expoentes) e subscritos (índices), respectivamente.  
		Exemplo: \verb|$x^2$| → $x^2$ e \verb|$a_i$| → $a_i$.
		
		\item \textbf{\{} e \textbf{\}} — delimitam grupos de comandos ou escopos.  
		Exemplo: \verb|\textbf{negrito}| → \textbf{negrito}.
		
		\item \textbf{\~{}} — gera o caractere til. Também pode ser usado para criar um espaço inquebrável, como em \verb|João~Silva|, garantindo que o nome não seja dividido no fim da linha.
		
		\item \textbf{\textbackslash} — introduz comandos no \LaTeX{}, como \verb|\textbf|, \verb|\section| etc.
		
		\item \textbf{\&} — separa colunas dentro de tabelas e ambientes matemáticos, como \verb|tabular| ou \verb|align|.
	\end{itemize}
	
	\bigskip
	\textbf{Dica:} se precisar exibir muitos desses símbolos com frequência (por exemplo, em tutoriais de programação ou explicações sobre LaTeX), considere usar o ambiente \verb|verbatim| ou o comando \verb|\verb|, que mostram o texto exatamente como digitado — sem que o \LaTeX{} tente interpretá-lo como código.
	
	
	% =============================
	\section{Salvando como PDF}
	
	Depois de escrever seu documento em \LaTeX{}, o passo final é gerar o arquivo em formato \texttt{.pdf}. Esse é o formato mais usado para compartilhamento e impressão de trabalhos, artigos e relatórios. O \LaTeX{} transforma automaticamente o código-fonte (\texttt{.tex}) em um PDF formatado com alta qualidade tipográfica.
	
	\bigskip
	\subsection*{Usando o Overleaf}
	
	Se você estiver utilizando o \textbf{Overleaf} (editor online), o processo é totalmente automático.  
	A cada alteração feita no código, o Overleaf recompila o projeto e exibe o PDF atualizado no painel à direita.  
	
	\begin{itemize}
		\item O botão \textbf{“Recompile”} (ou o atalho \texttt{Ctrl + Enter}) força uma nova compilação manual, caso o PDF não atualize automaticamente.
		\item O arquivo PDF final pode ser baixado clicando em \textbf{“Download PDF”} no menu superior.
		\item O Overleaf utiliza por padrão o compilador \verb|pdfLaTeX|, mas você também pode mudar para \verb|XeLaTeX| ou \verb|LuaLaTeX| nas configurações do projeto, se precisar de fontes especiais ou idiomas diferentes.
	\end{itemize}
	
	\bigskip
	\subsection*{Usando um editor local (TeXstudio, TeXworks, etc.)}
	
	Se preferir editar e compilar localmente no seu computador, você precisará de uma distribuição \LaTeX{} instalada — como o \textbf{TeX Live} (Linux e Windows) ou o \textbf{MiKTeX} (Windows).  
	
	Depois de instalada, basta abrir o seu arquivo \texttt{.tex} em um editor como o \textbf{TeXstudio} e seguir um dos métodos abaixo:
	
	\begin{itemize}
		\item Clique no botão \textquotedblleft \textbf{Reproduzir} \textquotedblright (geralmente representado por um triângulo verde) para compilar e gerar o PDF automaticamente;
		\item Ou, se preferir usar o terminal, digite:
		\begin{verbatim}
			pdflatex arquivo.tex
		\end{verbatim}
		Isso cria o arquivo \texttt{arquivo.pdf} na mesma pasta onde está o código-fonte.
	\end{itemize}
	
	\bigskip
	\textbf{Dica:} se o seu documento contém referências, sumário ou citações, é recomendável compilar duas vezes seguidas. A primeira compilação coleta as informações, e a segunda atualiza os números de páginas e referências cruzadas.
	
	\bigskip
	\subsection*{Erros comuns e soluções}
	
	Durante a compilação, podem ocorrer erros simples, como:
	\begin{itemize}
		\item Arquivo de imagem não encontrado — verifique o caminho e a extensão;
		\item Falta de um comando \verb|\end{...}| — revise se todos os ambientes estão corretamente fechados;
	\item Problemas com acentuação — se aparecerem caracteres estranhos, mude o compilador para \verb|XeLaTeX| e adicione o pacote:
	\begin{verbatim}
		\usepackage[utf8]{inputenc}
	\end{verbatim}
	\end{itemize}
	
	\bigskip
	\textbf{Resumo:}
	\begin{itemize}
	\item No \textbf{Overleaf}: o PDF é gerado automaticamente.
	\item No \textbf{computador}: compile manualmente com o TeXstudio ou pelo terminal usando \verb|pdflatex|.
	\end{itemize}
	
	Esses procedimentos garantem que seu documento final seja exportado com qualidade profissional, pronto para apresentação, publicação ou impressão.

	
\end{document}
